A dimensão $1/T$ da eq. (3.16) é confirmada pelo mesmo procedimento dimensional aplicado ao problema anterior.

\textbf{Tabela de dimensões das grandezas envolvidas:}
\begin{itemize}
\item Comprimento característico $L$: $[L] = L$
\item Aceleração gravitacional $g$: $[g] = L\,T^{-2}$
\item Velocidade $v$: $[v] = L\,T^{-1}$
\item Força $F$: $[F] = M\,L\,T^{-2}$
\item Massa $m$: $[m] = M$
\end{itemize}

\textbf{Verificação dimensional:}
Substituindo cada variável da eq. (3.16) por suas dimensões e realizando as operações algébricas sobre os expoentes, as potências de $M$ e $L$ se cancelam mutuamente, restando somente $T^{-1}$:
\[
[\text{eq.}\ (3.16)]\;=\; M^0 \cdot L^0 \cdot T^{-1} = T^{-1}.\quad\checkmark
\]
Portanto, a dimensão da eq. (3.16) é $\boxed{1/T}$, consistente com uma frequência ou taxa de variação.

